% META: {"id": "1NSI-RES-ME-004", "chapitre": "1NSI-RESEAUX", "type_objet": "methode", "capacites": ["P-ARCH-04A", "P-ARCH-04B"], "capacites_codes": ["C4", "C5"], "methodes": ["M4"], "mode_creation": "ex_nihilo", "sources_inspiration": [], "status": "needs_review"}
% BEGIN-VERIFY
% etat = {"ventilation": False}
% def decider_ventilation(co2, seuil_haut=1000, seuil_bas=800):
%     if etat["ventilation"]:
%         if co2 < seuil_bas:
%             etat["ventilation"] = False
%     else:
%         if co2 > seuil_haut:
%             etat["ventilation"] = True
%     return etat["ventilation"]
% mesures = [900, 1000, 1001, 850, 800, 799]
% assert [decider_ventilation(c) for c in mesures] == [False, False, True, True, True, False]
% END-VERIFY
\begin{fichemethode}{M4}{Du cahier des charges au programme : capteur, décision, actionneur --- puis tester les cas limites}

\textbf{Quand l'utiliser ?} Quand l'énoncé décrit un système physique (thermostat,
arrosage, éclairage, ventilation) et demande d'identifier ses composants, d'écrire la
fonction qui décide, ou de vérifier qu'elle respecte le cahier des charges.

\textbf{Pas à pas --- identifier les rôles :}
\begin{enumerate}
  \item \textbf{Le capteur est ce qui mesure} : il transforme une grandeur physique en un
        nombre que le programme lit. Pose la question « qu'est-ce qui produit la valeur en
        entrée ? ».
  \item \textbf{L'actionneur est ce qui agit} : il reçoit un ordre du programme et modifie
        le monde. Pose la question « qu'est-ce qui change quand le programme décide ? ».
        Un relais, un moteur, une pompe, une lampe agissent ; un thermomètre, une
        photorésistance, un détecteur mesurent.
\end{enumerate}

\textbf{Pas à pas --- programmer la décision :}
\begin{enumerate}
  \item \textbf{Traduis chaque règle du cahier des charges par une comparaison}, en
        respectant le mot exact : « sous » ou « au-dessus » se traduisent par \lstinline{<}
        ou \lstinline{>} stricts ; « à partir de » par \lstinline{<=} ou \lstinline{>=}.
  \item \textbf{Cherche si la décision dépend de l'état précédent.} Deux seuils (un pour
        démarrer, un pour arrêter) signifient que le programme doit \emph{mémoriser} s'il est
        en marche : la fonction consulte un état, puis le met à jour.
  \item \textbf{Dresse un tableau de tests} avant de conclure : une mesure franchement en
        dessous, une franchement au-dessus, \textbf{chaque seuil exactement}, et une
        séquence qui traverse les seuils dans les deux sens.
\end{enumerate}

\textbf{Exemple rédigé.} Cahier des charges d'une salle de classe : la ventilation
\emph{démarre} quand le CO$_2$ dépasse $1000$ ppm et ne s'\emph{arrête} que lorsqu'il
redescend sous $800$ ppm. Le capteur est la sonde de CO$_2$ ; l'actionneur est le
ventilateur. Deux seuils : il faut un état \lstinline{etat["ventilation"]}. Si la
ventilation tourne, on ne teste que l'arrêt (\lstinline{co2 < 800}) ; sinon, que le
démarrage (\lstinline{co2 > 1000}). Tableau de tests sur la séquence
$900, 1000, 1001, 850, 800, 799$ : arrêt, arrêt ($1000$ n'est pas \emph{au-dessus} de
$1000$), marche, marche, marche ($800$ n'est pas \emph{sous} $800$), arrêt. Soit
\lstinline{[False, False, True, True, True, False]}.

\textbf{Erreur fréquente.} Confondre les rôles : un capteur n'agit sur rien, un actionneur ne
mesure rien. Ou n'écrire qu'un seul seuil quand le cahier des charges en donne deux : la
ventilation clignoterait autour de $1000$ ppm. Ou ne tester que le scénario de
démonstration : c'est aux seuils exacts que le programme se trompe.

\end{fichemethode}
